\scanpage{196}
% manuscript position: top
\mspage{112}
% manuscript position: 1
la démonstration plus \unsure{sérieusement|sérieusemt}.

% manuscript position: 2
Soit donc $F$ \deleted{$\unsure{\in\mathrm{Filt}^\pi|=\mathrm{Filt}^\pi}$} \inserted{$\in\operatorname{Ob} M$}, \deleted{\unsure{considérons}} considérons
% manuscript position: 3
\[ J=C_{/F}\,,\quad \text{\deleted{($\unsure{\operatorname{card}\mathrm{Fl}\,J\leq\pi|{}\text{de }\operatorname{card}\mathrm{Fl}\,J\leq\pi}$)}}, \]
% manuscript position: 4–5
c'est une catégorie stable par \deleted{\unsure{sommes directes}} \inserted{$\varinjlim$ de cardinal} $\leq\pi_0$, donc grande devant $\pi_0$, et on a
% manuscript position: 6
\[ F\simeq\varinjlim_{j\in J}x_j\quad(\varinjlim\text{ dans }M) \]
% manuscript position: 7–8
en désignant par $j\longmapsto x_j$ le foncteur canonique $C_{/F}\longrightarrow C$.

% manuscript position: 9–12
Soit donc $I$ l'ens des sous-catégories \deleted{\unsure{pleines}}\inserted{$J'$} de $J$ qui sont \uline{de cardinal} $\leq\pi$ \uline{et} \uline{grandes devant} $\pi_0$, $I$ ordonné par inclusion, et pour tout $J'\in I$, posons
% manuscript position: 13
\[ X_{J'}=\varinjlim_{j\in J'}x_j\,. \]
% manuscript position: 14–16
\deleted{On a alors} \deleted{Alors \unsure{$\forall$} $J'\longmapsto X_{J'}$ est un foncteur $I\longrightarrow\mathrm{Filt}$} \deleted{\unsure{Cette|Celle}}

% manuscript position: 16–20
Comme $J'$ est grand devant $\pi_0$, la limite inductive dans $M$ est aussi une limite dans $\widehat C$, d'où résulte \deleted{que} aussitôt, comme $\operatorname{card}\,\text{\deleted{$\unsure{F|J}$}}\,\operatorname{Ob} J'\mathrel{\unsure{\subset|<}}\operatorname{card}\mathrm{Fl}\,J'\leq\pi$, que $X_{J'}$ est dans $\mathrm{Filt}^\pi$. Alors
% manuscript position: 21
\[ J'\longmapsto X_{J'}\in\operatorname{Ob}\mathrm{Filt}^{\unsure{\pi|n}} \]
% manuscript position: 22
est un foncteur
% manuscript position: 23
\[ I\longrightarrow\mathrm{Filt}^\pi\,, \]
% manuscript position: 24–27
dont la limite inductive est celle des $x_j$, donc isom. \unsure{;|à} $F$. D'autre part il est immédiat que $I$ est grd devant $\pi$, i.e. que si $(J'_\alpha)_{\alpha\in A}$ \deleted{est}
% manuscript position: foot
% archive sheet number: 195
